\subsection{二维加权 \texorpdfstring{$\zeta(z;u,v)$}{zeta(z;u,v)} 的 atom/core 相图、联合 residue shell 与奇偶迹律}\label{subsec:conclusion-realinput40-uv-atomcore-phase-shell-trace}
上一小节把长度 \(2\) 单原子在一参数输出位势下的 primitive surgery、residue 偏移与 Abel 常数差额压缩为完全显式的有限维公式。对附录 \ref{app:real-input-40-zeta-uv} 中新引入的二维联合势
\[
\Delta(z;u,v)=\det(I-zM_{u,v})=(vz^2-1)P_8(z;u,v)
\]
而言，同一单原子机制还可进一步提升为关于主谱机制切换、二维 Abel 修正、collision--output 联合 residue 壳以及 atom-dominant 区域奇偶迹律的四条终端命题。

\begin{theorem}[二维加权谱半径的 atom/core 精确分裂]\label{thm:conclusion-realinput40-uv-atomcore-spectral-radius-splitting}
沿用附录 \ref{app:real-input-40-zeta-uv} 的记号。对任意 \(u\in\CC\) 与 \(v>0\)，定义
\[
\rho_{\mathrm{full}}(u,v):=\rho(M_{u,v}),
\qquad
\rho_{\mathrm{core}}(u,v):=\max\bigl\{|\lambda|:\ P_8(\lambda^{-1};u,v)=0\bigr\}.
\]
则有严格恒等式
\[
\boxed{\;
\rho_{\mathrm{full}}(u,v)=\max\bigl\{\sqrt v,\rho_{\mathrm{core}}(u,v)\bigr\}.\;}
\]
进一步定义分界集合
\[
\mathcal C:=\bigl\{(u,v)\in\CC\times\RR_{>0}:\ \rho_{\mathrm{core}}(u,v)=\sqrt v\bigr\}.
\]
则：
\begin{enumerate}
  \item 在 \emph{core-dominant} 区域 \(\rho_{\mathrm{core}}(u,v)>\sqrt v\) 上，最大模谱完全由 \(P_8\) 的根给出；
  \item 在 \emph{atom-dominant} 区域 \(\rho_{\mathrm{core}}(u,v)<\sqrt v\) 上，最大模谱完全由显式谱支 \(\pm\sqrt v\) 给出；
  \item 分界集合 \(\mathcal C\) 是主导机制发生切换的唯一可能位置。
\end{enumerate}
\leanverified{spectral\_radius\_max\_splitting}
\leanverified{core\_dominant\_max}
\leanverified{atom\_dominant\_max}
\leanverified{critical\_locus\_max}
\leanverified{paper\_conclusion\_realinput40\_uv\_atomcore\_seeds}
\leanverified{paper\_conclusion\_realinput40\_uv\_atomcore\_package}
\end{theorem}
\begin{proof}
由定理 \ref{thm:real-input-40-det-uv-2plus8}，
\[
\Delta(z;u,v)=(vz^2-1)P_8(z;u,v).
\]
因此 \(M_{u,v}\) 的全部非零特征值，恰是方程
\[
vz^2-1=0
\qquad\text{与}\qquad
P_8(z;u,v)=0
\]
之零点倒数的并。再由推论 \ref{cor:real-input-40-zeta-uv-sqrtv-eigs}，前一因子始终给出两条显式特征值
\[
\lambda_\pm=\pm\sqrt v.
\]
其余全部非零谱则完全由 \(P_8\) 承担。故谱半径只能取自这两部分中的较大者，从而
\[
\rho_{\mathrm{full}}(u,v)=\max\bigl\{\sqrt v,\rho_{\mathrm{core}}(u,v)\bigr\}.
\]
后三条只是这一 \(\max\)-结构的直接分区解释。证毕。
\end{proof}

\begin{theorem}[二维 primitive 手术的可加性与 Abel 常数修正]\label{thm:conclusion-realinput40-uv-primitive-surgery-abel-shift}
\leanverified{paper\_conclusion\_realinput40\_uv\_primitive\_surgery\_abel\_shift}
沿用附录 \ref{app:real-input-40-zeta-uv} 的记号，并写
\[
\Delta(z;u,v)=(1-vz^2)\,\mathcal Q_8(z;u,v),
\qquad
\mathcal Q_8(z;u,v):=-P_8(z;u,v),
\]
\[
\zeta(z;u,v):=\Delta(z;u,v)^{-1},
\qquad
\zeta_{\mathrm{core}}(z;u,v):=\mathcal Q_8(z;u,v)^{-1}.
\]
\[
P(z;u,v)=\sum_{n\ge 1}p_n(u,v)z^n,
\qquad
P_{\mathrm{core}}(z;u,v)=\sum_{n\ge 1}p_n^{\mathrm{core}}(u,v)z^n.
\]
则始终有
\[
\boxed{\;
P(z;u,v)=vz^2+P_{\mathrm{core}}(z;u,v).\;}
\]
若进一步 \(z_\star(u,v)\) 是 \(\zeta(z;u,v)\) 的主极点且由 core 因子给出，即
\[
\rho_{\mathrm{core}}(u,v)>\sqrt v
\qquad
\bigl(\text{从而 }1-vz_\star(u,v)^2\neq 0\bigr),
\]
并记 \(\log\mathfrak M(u,v)\)、\(\log\mathfrak M_{\mathrm{core}}(u,v)\) 为以 \(z_\star(u,v)\) 归一化的 Abel 有限部常数，则有严格公式
\[
\boxed{\;
\log\mathfrak M(u,v)=\log\mathfrak M_{\mathrm{core}}(u,v)+v\,z_\star(u,v)^2.\;}
\]
\leanverified{abel\_constant\_shift}
\end{theorem}
\begin{proof}
固定碰撞参数 \(u\)，把 \(v\) 视为单一 Euler 权变量。由
\[
\zeta(z;u,v)=\frac{1}{1-vz^2}\,\zeta_{\mathrm{core}}(z;u,v).
\]
定理 \ref{thm:xi-finite-euler-primitive-surgery-additivity} 在单原子情形
\[
(m,\ell,E)=(1,2,1)
\]
下立即给出
\[
P(z;u,v)=vz^2+P_{\mathrm{core}}(z;u,v).
\]
若再假设 \(z_\star(u,v)\) 来自 core 而非原子因子，则
\[
1-vz_\star(u,v)^2\neq 0.
\]
于是可将定理 \ref{thm:xi-atomic-witt-surgery-pressure-ldp-abel-invariance} 逐字应用于同一单原子情形
\[
(m,\ell,E)=(1,2,1),
\]
其中活动权变量由原定理中的 \(u\) 改写为此处的 \(v\)，而碰撞参数 \(u\) 仅作为旁参数保留。由该定理立得
\[
\log\mathfrak M(u,v)=\log\mathfrak M_{\mathrm{core}}(u,v)+v\,z_\star(u,v)^2.
\]
证毕。
\end{proof}

\begin{theorem}[二维 collision--output residue 计数的单原子 delta 壳]\label{thm:conclusion-realinput40-uv-collision-output-delta-shell}
固定整数 \(m\ge 2\)。对 \(n\ge 1\)、\(c\ge 0\) 与 \(r\in\ZZ/m\ZZ\)，令
\[
A_{m,r}^{\mathrm{full}}(n;c)
\]
表示长度为 \(n\) 的 periodic 对象中，碰撞次数恰为 \(c\)、且输出 \(1\) 的次数模 \(m\) 等于 \(r\) 的计数；令
\[
P_{m,r}^{\mathrm{full}}(n;c)
\]
表示相应的 primitive 计数。对 core 部分同样定义
\[
A_{m,r}^{\mathrm{core}}(n;c),\qquad
P_{m,r}^{\mathrm{core}}(n;c).
\]
则有精确差分公式
\[
\boxed{\;
A_{m,r}^{\mathrm{full}}(n;c)-A_{m,r}^{\mathrm{core}}(n;c)
=
2\,\mathbf 1_{\{2\mid n\}}
\mathbf 1_{\{c=0\}}
\mathbf 1_{\{r\equiv n/2\;(\mathrm{mod}\,m)\}},\;}
\]
\[
\boxed{\;
P_{m,r}^{\mathrm{full}}(n;c)-P_{m,r}^{\mathrm{core}}(n;c)
=
\mathbf 1_{\{n=2\}}
\mathbf 1_{\{c=0\}}
\mathbf 1_{\{r\equiv 1\;(\mathrm{mod}\,m)\}}.\;}
\]
因此，二维计数世界中的单原子贡献在 primitive 层完全集中于唯一 joint atom
\[
(|\tau|,\mathrm{collision},\mathrm{output})=(2,0,1),
\]
而在 periodic 层则仅沿
\[
r\equiv \frac{n}{2}\pmod m
\]
这条 residue 壳传播。
\leanverified{delta\_shell\_indicator}
\end{theorem}
\begin{proof}
由定理 \ref{thm:conclusion-realinput40-uv-primitive-surgery-abel-shift}，
\[
P(z;u,v)-P_{\mathrm{core}}(z;u,v)=vz^2.
\]
这说明原子 primitive 项只含唯一单项式 \(u^0v^1z^2\)：长度恰为 \(2\)，碰撞次数恰为 \(0\)，输出次数恰为 \(1\)。按输出次数模 \(m\) 分组，立即得到
\[
P_{m,r}^{\mathrm{full}}(n;c)-P_{m,r}^{\mathrm{core}}(n;c)
=
\mathbf 1_{\{n=2\}}
\mathbf 1_{\{c=0\}}
\mathbf 1_{\{r\equiv 1\;(\mathrm{mod}\,m)\}}.
\]

另一方面，periodic/ghost 层的原子部分满足
\[
\log\frac{1}{1-vz^2}
=
\sum_{q\ge 1}\frac{v^q}{q}z^{2q},
\]
故对应 ghost 系数差为
\[
a_n^{\mathrm{full}}(u,v)-a_n^{\mathrm{core}}(u,v)
=
2v^{n/2}\mathbf 1_{\{2\mid n\}}.
\]
按本文的 Witt--ghost 记号，这正是长度 \(n\) periodic 计数的原子生成多项式。由于该差分不含 \(u\)，碰撞次数只能是 \(c=0\)；而单项式 \(v^{n/2}\) 则表明输出次数恰为 \(n/2\)。再按模 \(m\) 归类，便得到
\[
A_{m,r}^{\mathrm{full}}(n;c)-A_{m,r}^{\mathrm{core}}(n;c)
=
2\,\mathbf 1_{\{2\mid n\}}
\mathbf 1_{\{c=0\}}
\mathbf 1_{\{r\equiv n/2\;(\mathrm{mod}\,m)\}}.
\]
证毕。
\end{proof}

\begin{theorem}[二维 atom-dominant 区域的奇偶迹律与次谱缺口]\label{thm:conclusion-realinput40-uv-atom-dominant-trace-law}
定义 atom-dominant 区域
\[
\mathcal A:=
\bigl\{(u,v)\in\CC\times\RR_{>0}:\ \rho_{\mathrm{core}}(u,v)<\sqrt v\bigr\},
\]
其中 \(\rho_{\mathrm{core}}(u,v)\) 由定理 \ref{thm:conclusion-realinput40-uv-atomcore-spectral-radius-splitting} 给出。则对任意 \((u,v)\in\mathcal A\)，存在常数 \(\eta(u,v)\in(0,1)\) 使得对所有 \(n\ge 1\) 都有
\[
\operatorname{Tr}(M_{u,v}^n)=
\begin{cases}
2v^{n/2}+O\!\bigl((\eta(u,v)\sqrt v)^n\bigr),& 2\mid n,\\[4pt]
O\!\bigl((\eta(u,v)\sqrt v)^n\bigr),& 2\nmid n.
\end{cases}
\]
并且若定义次谱半径
\[
\Lambda_2(u,v):=
\max\bigl\{|\lambda|:\ \lambda\in\Spec(M_{u,v})\setminus\{\pm\sqrt v\}\bigr\},
\]
则有严格缺口
\[
\boxed{\;
1-\frac{\Lambda_2(u,v)}{\sqrt v}>0.\;}
\]
\leanverified{paper\_conclusion\_realinput40\_uv\_atom\_dominant\_trace\_law}
\leanverified{even\_sign\_sum}
\leanverified{odd\_sign\_sum}
\leanverified{spectral\_gap\_pos}
\end{theorem}
\begin{proof}
在 \((u,v)\in\mathcal A\) 上，由定理 \ref{thm:conclusion-realinput40-uv-atomcore-spectral-radius-splitting}，
\[
\rho(M_{u,v})=\sqrt v,
\qquad
\rho_{\mathrm{core}}(u,v)<\sqrt v.
\]
再由推论 \ref{cor:real-input-40-zeta-uv-sqrtv-eigs}，两条显式 rigid branches 的迹贡献恰为
\[
(\sqrt v)^n+(-\sqrt v)^n
=
\begin{cases}
2v^{n/2},&2\mid n,\\
0,&2\nmid n.
\end{cases}
\]
其余全部特征值都属于 core 部分，模长至多为 \(\rho_{\mathrm{core}}(u,v)\)。任选
\[
\eta(u,v)\in\left(\frac{\rho_{\mathrm{core}}(u,v)}{\sqrt v},1\right),
\]
则所有其余 Jordan 块的贡献都可统一吸收到
\[
O\!\bigl((\eta(u,v)\sqrt v)^n\bigr)
\]
中，因为固定次数的多项式因子可由更大的指数底吸收。于是得到所示奇偶迹律。

最后，由定义
\[
\Lambda_2(u,v)\le \rho_{\mathrm{core}}(u,v)<\sqrt v,
\]
故
\[
1-\frac{\Lambda_2(u,v)}{\sqrt v}>0.
\]
证毕。
\end{proof}

\begin{theorem}[atom-dominant 区域的主谱、奇偶迹与 residue 壳三签名同步]\label{thm:conclusion-realinput40-uv-atom-dominant-three-signature-synchrony}
\leanverified{paper\_conclusion\_realinput40\_uv\_atom\_dominant\_three\_signature\_synchrony}
在 atom-dominant 区域
\[
\mathcal A
=
\bigl\{(u,v)\in\CC\times\RR_{>0}:\ \rho_{\mathrm{core}}(u,v)<\sqrt v\bigr\}
\]
内，单 balanced atom \((1-vz^2)^{-1}\) 同时强制三种彼此独立而又同步的可见签名：
\begin{enumerate}
  \item 主谱签名
  \[
  \rho_{\mathrm{full}}(u,v)=\sqrt v;
  \]
  \item 奇偶迹签名：存在 \(\eta(u,v)\in(0,1)\) 使
  \[
  \operatorname{Tr}(M_{u,v}^n)=
  \begin{cases}
  2v^{n/2}+O\!\bigl((\eta(u,v)\sqrt v)^n\bigr),& 2\mid n,\\[4pt]
  O\!\bigl((\eta(u,v)\sqrt v)^n\bigr),& 2\nmid n;
  \end{cases}
  \]
  \item collision--output residue 壳签名
  \[
  A_{m,r}^{\mathrm{full}}(n;c)-A_{m,r}^{\mathrm{core}}(n;c)
  =
  2\,\mathbf 1_{\{2\mid n\}}
  \mathbf 1_{\{c=0\}}
  \mathbf 1_{\{r\equiv n/2\;(\mathrm{mod}\,m)\}},
  \]
  而 primitive 层仅保留唯一的 joint atom
  \[
  P_{m,r}^{\mathrm{full}}(n;c)-P_{m,r}^{\mathrm{core}}(n;c)
  =
  \mathbf 1_{\{n=2\}}
  \mathbf 1_{\{c=0\}}
  \mathbf 1_{\{r\equiv 1\;(\mathrm{mod}\,m)\}}.
  \]
\end{enumerate}
因此，在 \(\mathcal A\) 内，单个 balanced atom 不仅支配谱半径，而且同时支配迹的偶奇律与 collision--output residue 壳。
\end{theorem}
\begin{proof}
主谱签名正是定理 \ref{thm:conclusion-realinput40-uv-atomcore-spectral-radius-splitting} 在 \(\mathcal A\) 上的直接特例。奇偶迹签名由定理 \ref{thm:conclusion-realinput40-uv-atom-dominant-trace-law} 给出。至于 residue 壳与 primitive joint atom，则是定理 \ref{thm:conclusion-realinput40-uv-collision-output-delta-shell} 的精确内容。将三者并置，即得所述同步律。证毕。
\end{proof}

上述五条命题表明，二维联合势中的单原子因子 \((1-vz^2)^{-1}\) 已不再只是对 core 的低阶修饰，而是一个能够与八次 core 谱直接竞争的独立主导机制：其主谱侧表现为 atom/core 的严格相图分裂，常数层表现为完全可加的二维 Abel 位移，而在 joint collision--output 计数侧则退化为唯一的 delta 型原子及其 periodic residue 壳；一旦进入 atom-dominant 区域，这三种可见量又被进一步锁成同一条主谱--奇偶迹--residue 壳同步律。因而二维谱、primitive surgery、有限部常数与 residue 读出便在同一组显式公式中完成闭合。

\leanverified{paper\_conclusion\_realinput40\_root\_unity\_ghost\_complete\_primitive\_degenerate}
\begin{theorem}[根值单位层析仅在 ghost 扇区完备]\label{thm:conclusion-realinput40-root-unity-ghost-complete-primitive-degenerate}
设可剥离原子 Witt 因子为
\[
\zeta_{\mathrm{atom}}(z,u)=(1-uz^2)^{-1},
\]
并记
\[
\log \zeta_{\mathrm{atom}}(z,u)=\sum_{n\ge 1}\frac{a_n^{\mathrm{atom}}(u)}{n}z^n,
\qquad
P_{\mathrm{atom}}(z,u)=\sum_{n\ge 1}p_n^{\mathrm{atom}}(u)z^n.
\]
则有严格闭式
\[
a_n^{\mathrm{atom}}(u)=2u^{n/2}\mathbf 1_{\{2\mid n\}},
\qquad
p_n^{\mathrm{atom}}(u)=u\,\mathbf 1_{\{n=2\}}.
\]
对任意整数 \(m\ge 2\) 与 \(\omega_m=e^{2\pi i/m}\)，进一步有精确平均律
\[
\frac{1}{m}\sum_{j=0}^{m-1}a_n^{\mathrm{atom}}(\omega_m^j)
=
2\,\mathbf 1_{\{2m\mid n\}}.
\]
因此，全部根值单位探针族 \(\{\omega_m\}_{m\ge 2}\) 在 ghost 层给出严格的 Dirichlet 梳齿层析；而 primitive 层除唯一长度 \(2\) 的原子脉冲外不生成任何新支撑。换言之，一切仅依赖有限根值单位采样的 \(\pi\)-通道，至多恢复偶长度的整除几何，而不能恢复 primitive 轨道结构。
\end{theorem}
\begin{proof}
由
\[
\log(1-uz^2)^{-1}
=
\sum_{q\ge 1}\frac{u^q}{q}z^{2q},
\]
直接读得
\[
a_n^{\mathrm{atom}}(u)=2u^{n/2}\mathbf 1_{\{2\mid n\}}.
\]
另一方面，原子 primitive 项只有唯一长度 \(2\) 的单项式，故
\[
P_{\mathrm{atom}}(z,u)=uz^2,
\qquad
p_n^{\mathrm{atom}}(u)=u\,\mathbf 1_{\{n=2\}}.
\]

对单位根平均，若 \(n\) 为奇数则左端显然为零；若 \(n=2q\) 为偶数，则
\[
\frac{1}{m}\sum_{j=0}^{m-1}a_n^{\mathrm{atom}}(\omega_m^j)
=
\frac{2}{m}\sum_{j=0}^{m-1}\omega_m^{jq}
=
2\,\mathbf 1_{\{m\mid q\}}
=
2\,\mathbf 1_{\{2m\mid n\}}.
\]
故 ghost 层平均恰形成整除条件 \(2m\mid n\) 的梳齿，而 primitive 层始终只保留单一长度 \(2\) 的原子。证毕。
\end{proof}

\leanverified{paper\_conclusion\_periodic\_dirichlet\_shadow\_hurwitz\_finite\_closure}
\begin{theorem}[周期 Dirichlet 阴影的 Hurwitz 有限闭包]\label{thm:conclusion-periodic-dirichlet-shadow-hurwitz-finite-closure}
设 \(\kappa_A:\NN\to\CC\) 为周期 \(d_m\) 的核函数，并定义
\[
Z_A(s):=\sum_{b\ge 1}\frac{\kappa_A(b)}{b^s},
\qquad \Re s>1.
\]
则有精确 Hurwitz 分解
\[
Z_A(s)=d_m^{-s}\sum_{r=1}^{d_m}\kappa_A(r)\,
\zeta\!\left(s,\frac{r}{d_m}\right),
\]
且在 \(s=1\) 处仅有唯一简单极点，并满足
\[
\operatorname*{Res}_{s=1}Z_A(s)
=
\mu_A
:=
\frac{1}{d_m}\sum_{r=1}^{d_m}\kappa_A(r).
\]
从而，每一个有限周期的标量 Dirichlet 阴影，在其表示层面上都完全等价于一个有限残数账本
\[
\bigl(d_m;\kappa_A(1),\dots,\kappa_A(d_m)\bigr),
\]
其解析内容在标量层不再携带额外自由度。
\end{theorem}
\begin{proof}
由周期性 \(\kappa_A(b+d_m)=\kappa_A(b)\)，分块求和得
\[
Z_A(s)
=
\sum_{r=1}^{d_m}\sum_{q\ge 0}\frac{\kappa_A(qd_m+r)}{(qd_m+r)^s}
=
\sum_{r=1}^{d_m}\kappa_A(r)\sum_{q\ge 0}(qd_m+r)^{-s}.
\]
提出公因子 \(d_m^{-s}\) 即得
\[
Z_A(s)
=
d_m^{-s}\sum_{r=1}^{d_m}\kappa_A(r)\sum_{q\ge 0}\left(q+\frac{r}{d_m}\right)^{-s}
=
d_m^{-s}\sum_{r=1}^{d_m}\kappa_A(r)\,
\zeta\!\left(s,\frac{r}{d_m}\right).
\]
又因每一项 \(\zeta(s,a)\) 在 \(s=1\) 的留数均为 \(1\)，故
\[
\operatorname*{Res}_{s=1}Z_A(s)
=
\frac{1}{d_m}\sum_{r=1}^{d_m}\kappa_A(r)
=
\mu_A.
\]
这说明周期 Dirichlet 阴影在标量层面已完全坍缩为有限周期向量及其平均值。证毕。
\end{proof}

\begin{proposition}[素数根值单位探针仅具二次分辨率]\label{prop:conclusion-prime-root-unity-probe-quadratic-resolution}
\leanverified{paper\_conclusion\_prime\_root\_unity\_probe\_quadratic\_resolution}
令 \(u=e^{it}\)，\(\rho(t):=\rho(e^{it})\)，并记 \(\lambda=\rho(1)=\varphi^2\)。则当 \(t\to 0\) 时有展开
\[
\log\frac{\rho(t)}{\lambda}
=
-\frac{6\sqrt5-5}{250}\,t^2
+\frac{7+24\sqrt5}{75000}\,t^4
+O(t^6).
\]
特别地，对奇素数 \(p\) 的最小非平凡根值单位角
\[
t_p:=\frac{2\pi}{p},
\]
有
\[
\log\frac{\rho(t_p)}{\lambda}
=
-\frac{6\sqrt5-5}{250}\left(\frac{2\pi}{p}\right)^2
+\frac{7+24\sqrt5}{75000}\left(\frac{2\pi}{p}\right)^4
+O(p^{-6}).
\]
因此，素数模根值单位探针对谱宿主的最小可分辨偏移只有 \(p^{-2}\) 级；结合定理 \ref{thm:conclusion-realinput40-root-unity-ghost-complete-primitive-degenerate}，这类探针本质上仍停留在 ghost 整除层析侧，而不能逼近宿主的 sheet 几何。
\end{proposition}
\begin{proof}
首个展开式即文稿中单位根扭曲局部展开的闭式结论。将 \(t=t_p=2\pi/p\) 直接代入，便得第二式。其主项为严格负的二次量，故根值单位探针的最小可分辨偏移量级只能达到 \(p^{-2}\)。而定理 \ref{thm:conclusion-realinput40-root-unity-ghost-complete-primitive-degenerate} 已表明有限根值单位平均在结构上只给出 ghost 层的整除梳齿，故这种二次分辨率提升并不触及单值化片结构。证毕。
\end{proof}


\begin{corollary}[根值单位层析的 ghost/primitive 秩分裂]\label{cor:conclusion-realinput40-root-unity-ghost-primitive-rank-split}
\leanverified{paper\_conclusion\_realinput40\_root\_unity\_ghost\_primitive\_rank\_split}
对 balanced atom
\[
\zeta_{\mathrm{atom}}(z,u)=(1-uz^2)^{-1},
\]
根值单位平均在 ghost 与 primitive 两层之间发生严格秩分裂：
\begin{enumerate}
  \item 在 ghost 层，对任意 \(m\ge 2\) 与 \(\omega_m=e^{2\pi i/m}\)，有
  \[
  \frac{1}{m}\sum_{j=0}^{m-1}a_n^{\mathrm{atom}}(\omega_m^j)
  =
  2\,\mathbf 1_{\{2m\mid n\}},
  \]
  因而它精确恢复偶长度可被哪些 \(m\) 整除的整除几何；
  \item 在 primitive 层，始终只有
  \[
  p_n^{\mathrm{atom}}(u)=u\,\mathbf 1_{\{n=2\}},
  \]
  因而全部根值单位采样在 primitive 方向上始终只有秩一分辨能力。
\end{enumerate}
换言之，根值单位层析对 ghost 扇区是完备的，对 primitive 扇区却退化为单一长度 \(2\) 的脉冲读出。
\end{corollary}
\begin{proof}
两条公式都已由定理 \ref{thm:conclusion-realinput40-root-unity-ghost-complete-primitive-degenerate} 给出。第一条说明 ghost 平均恰形成整除条件 \(2m\mid n\) 的 Dirichlet 梳齿，因此完整恢复偶长度的整除几何；第二条则表明 primitive 层无论考察多少个根值单位模数，都不会产生除长度 \(2\) 之外的新支撑。故秩分裂成立。证毕。
\end{proof}

进一步说，这个单原子机制在根值单位方向并不会把先前的三签名同步继续提升到 primitive 结构，而是发生一次严格的 ghost/primitive 秩分裂：periodic/ghost 层析被压缩成严格的 Dirichlet 梳齿，并把全部有限周期标量阴影锁死在 Hurwitz 残数账本中；primitive 层却始终只剩长度 \(2\) 的单脉冲。即使把素数模根值单位探针的角分辨率推进到 \(p^{-2}\) 级，它所细化的也仍只是 ghost 整除几何，而不是宿主的 sheet 结构。于是 atom/core 分裂、主谱--奇偶迹--residue 壳同步、ghost 梳齿与 Hurwitz 有限闭包便共同构成了二维 realinput40 核在 cyclotomic 方向上的完整可见边界。

\endinput
